Beaglehole et al. (\cite{beaglehole2026toward}) afirmam que os conceitos ``conhecidos'' pelos \textit{Large Language Models} (Grandes Modelos de Linguagem, em tradução livre) são, de fato, direções lineares no espaço de ativações do modelo, que tais direções podem ser extraídas por algoritmos de aprendizado de características como as \textit{Recursive Feature Machines} (Máquinas Recursivas de Características, em tradução livre, doravante RFM), e que, uma vez extraído um conceito, é possível aplicar \textbf{\textit{steering}} (direcionamento, em tradução livre) ao modelo, ou seja, empurrar suas saídas em direção ao conceito ou para longe dele por meio de computações puramente lineares, bem como realizar o \textbf{\textit{monitoring}} (monitoramento, em tradução livre) da presença do conceito pela leitura das ativações. Estas afirmações são estabelecidas sobre modelos da classe Llama, com 8 a 90 bilhões de parâmetros, e é portanto razoável perguntar se os mesmos resultados sobrevivem em escalas muito menores. Tais achados se sustentam em um modelo de apenas $\approx 9$ milhões de parâmetros? Note que um único \textit{checkpoint} não pode dizer como a maquinaria escala, uma vez que isso exigiria uma escada de tamanhos de modelo e constitui um experimento à parte, mas ele pode dizer se o mecanismo sobrevive por seus próprios méritos no piso da faixa, e uma resposta negativa ali limitaria toda a afirmação.

Este documento relata uma replicação controlada dos mecanismos supracitados, a saber, a extração via RFM, o \textit{steering} aditivo, o \textit{probing} por autovetores da AGOP e a comparação contra \textit{baselines} (métodos de referência) mais baratos, conduzida sobre o GuppyLM, um modelo de linguagem de $\approx 8{,}7$M de parâmetros. Deliberadamente abrimos mão dos resultados temáticos de Beaglehole et al., como o \textit{steering} ao longo de uma direção antirrecusa para expor a vulnerabilidade de um modelo a \textit{jailbreaks}, ou a subtração de uma direção de engano para que o modelo responda de forma mais honesta, já que o GuppyLM não carrega representações de tais matérias. Perguntamos, ao invés disso, se a extração linear de conceitos, o \textit{steering} aditivo e o \textit{probing} de ativações se sustentam em modelos de linguagem tão pequenos quanto 10M de parâmetros, restritos aos conceitos que o modelo é capaz de representar.

Em suma, a leitura dos conceitos transfere-se intacta, a escrita deles, isto é, o \textit{steering}, transfere-se como mecanismo e não como capacidade, e a comparação contra os \textit{baselines} é apenas parcialmente testável nesta escala.
